\chapter{Implementace}
\label{chap:implementace}

V této kapitole je detailně rozebrán samotný proces implementace navrženého řešení. Pozornost je věnována struktuře celého projektu a následně specifickým technickým výzvám spojeným s vývojem klientské aplikace v jazyce Go a úpravám systémové služby na platformě .NET.

\section{Struktura projektu a vývojové prostředí}
Jelikož se výsledný produkt skládá z několika oddělených, avšak spolupracujících částí, bylo zásadní správně navrhnout strukturu repozitáře. Repozitář je koncipován jako monorepo, které zastřešuje zdrojové kódy jak pro klientskou část (frontend), tak pro vrstvu systémové služby (backend). Tento přístup zjednodušuje správu závislostí a verzování.

Hlavní rozdělení repozitáře je následující:
\begin{itemize}
    \item \texttt{cli/} -- Obsahuje zdrojové kódy klientské aplikace napsané v jazyce Go. Kód je strukturován s ohledem na oddělení uživatelského rozhraní od komunikační logiky (adresáře \texttt{cmd/} a \texttt{internal/}).
    \item \texttt{services/} -- Zahrnuje .NET projekty představující systémovou službu, která běží na pozadí.
    \begin{itemize}
        \item \texttt{CLIService/} -- Nově vytvořený modul určený primárně pro obsluhu IPC komunikace s CLI klientem a řízení jeho požadavků.
        \item \texttt{GoodAccessService/} -- Původní jádro služby obstarávající infrastrukturu a samotné sestavování tunelů (např. OpenVPN, WireGuard).
    \end{itemize}
    \item \texttt{GoodAccess.Shared/} -- Sdílené doménové modely, rozhraní a utility využívané napříč různými .NET službami.
\end{itemize}

Architektura je tak striktně rozdělena na hloupého klienta (Go), který se stará pouze o prezentaci dat uživateli, a chytrou službu (.NET), kde probíhá veškerá složitá byznys logika.

\section{Implementace systémové služby (.NET)}
Aby bylo možné efektivně zpracovávat požadavky přicházející od CLI klienta, bylo nutné infrastrukturu na straně služby patřičně rozšířit. Vývoj probíhal iterativně počínaje zajištěním bezpečného úložiště a zprovozněním lokálního komunikačního kanálu.

\subsection{Meziprocesní komunikace (IPC)}
Základním kamenem architektury je nově vytvořený modul \texttt{CLIService}. Tento modul slouží jako most mezi uživatelskými vstupy a jádrem služby. Komunikace probíhá přes Unix Domain Sockets (UDS), což je na platformě Linux standardní a vysoce výkonný způsob výměny dat mezi lokálními procesy.

Služba vytvoří pojmenovanou rouru (named pipe), na které asynchronně naslouchá příchozím spojením od klienta. Požadavky jsou zasílány a vraceny ve formátu JSON, což usnadňuje deserializaci dat na obou stranách.

\begin{figure}[ht]
\begin{lstlisting}
// Ukázka inicializace IPC serveru pomocí NamedPipeServerStream pro Linux
var pipeServer = new NamedPipeServerStream(
    pipeName: "CoreFxPipe_ga-cli.sock",
    direction: PipeDirection.InOut,
    maxNumberOfServerInstances: 1,
    transmissionMode: PipeTransmissionMode.Byte,
    options: PipeOptions.Asynchronous);

await pipeServer.WaitForConnectionAsync(cancellationToken);
// Zpracování příchozích dat ze socketu...
\end{lstlisting}
\caption{Inicializace Unix Domain Socketu na straně .NET služby}
\label{fig:ipc_init}
\end{figure}

\subsection{Bezpečné úložiště (Secure Storage)}
Aplikace potřebuje uchovávat citlivá data jako jsou autentizační tokeny a privátní klíče. Původní implementace pro grafické rozhraní ukládala některá data do prostoru uživatele. Jelikož CLI je navrženo s možností trvalého spojení běžícího čistě na úrovni systému, muselo být úložiště přesunuto do bezpečného systémového adresáře.

Byla vytvořena třída reprezentující \textit{Secure Storage}. Fyzické uložení dat využívá soubory zabezpečené restriktivními přístupovými právy (\texttt{chmod 700}), čímž je zaručeno, že k souborům má přístup pouze vlastník procesu (v tomto případě uživatel \texttt{root}, pod kterým služba běží). Samotná data uvnitř souborů jsou navíc šifrována pomocí \texttt{DataProtection API}.

\subsection{Zpracování příkazů (CliMessenger)}
Za samotné odbavování příchozích zpráv zodpovídá komponenta \texttt{CliMessenger}. Každý požadavek ve formátu JSON obsahuje specifikaci příkazu (např. \texttt{login}, \texttt{connect}) a příslušné argumenty. \texttt{CliMessenger} tyto požadavky směruje na specifické služby v rámci .NET backendu.

V raných fázích vývoje, kdy ještě nebyla zcela implementována veškerá produkční logika a klientská aplikace se teprve formovala, vracela metoda \texttt{CliMessenger} provizorní pevně daná data (tzv. hardcoded odpovědi). Tento přístup umožnil paralelní vývoj -- UI se mohlo designovat a testovat proti konzistentnímu API bez čekání na dokončení komplexní aplikační logiky. Postupem času byly tyto provizorní odpovědi nahrazovány reálným napojením na síťové a autentizační subsystémy.

\section{Klientská aplikace a příkazy (Go)}
Na straně klienta bylo zvoleno využití jazyka Go, a to především pro jeho vynikající vlastnosti při vývoji CLI aplikací a schopnost kompilace do jediné statické binárky nezávislé na externích knihovnách na koncovém systému. 

\subsection{Struktura příkazů (Cobra)}
Základní aplikační kostra byla postavena na populárním frameworku Cobra. Tento nástroj umožňuje velmi snadno deklarovat stromovou strukturu příkazů, definovat argumenty (tzv. flags) a automaticky generovat nápovědu pro uživatele.

Hlavní příkaz \texttt{ga-cli} byl inicializován jako kořenový uzel (Root Command), ke kterému byly postupně přidávány dílčí příkazy odpovídající požadovaným akcím (login, setup, connect a další). Architektura frameworku Cobra navíc vybízí k přehlednému uspořádání kódu, kdy má každý příkaz svůj vlastní soubor v adresáři \texttt{cmd/}.

\begin{figure}[ht]
\begin{lstlisting}
// Ukázka inicializace kořenového příkazu
var rootCmd = &cobra.Command{
	Use:   "ga-cli",
	Short: "GoodAccess CLI klient",
	Long:  "Rozhraní příkazové řádky pro správu GoodAccess VPN na Linuxu.",
}

func init() {
	rootCmd.AddCommand(loginCmd)
	rootCmd.AddCommand(connectCmd)
	rootCmd.AddCommand(statusCmd)
	// Registrace dalších příkazů...
}
\end{lstlisting}
\caption{Definice struktury příkazů pomocí frameworku Cobra}
\label{fig:cobra_init}
\end{figure}

\subsection{Implementace příkazů a stavový report}
Jednotlivé příkazy byly do aplikace implementovány iterativním způsobem v logickém pořadí:
\begin{enumerate}
    \item \textbf{login / logout:} První krok zajišťující ověření uživatele a zisk patřičných tokenů k další komunikaci.
    \item \textbf{setup:} Konfigurace připojení – výběr vhodné brány (gateway) ze seznamu, nastavení požadovaného protokolu a definice parametru persistentní konektivity.
    \item \textbf{connect / disconnect:} Finální inicializace samotného šifrovaného tunelu, což původně zahrnovalo primárně integraci s \texttt{ProtoService} a stávajícím OpenVPN rozhraním.
\end{enumerate}

Během celého procesu implementace výše zmíněných příkazů se ukázal jako kritický podpůrný prvek příkaz \textbf{status}. Tento příkaz byl průběžně doplňován a využíván k rychlé diagnostice stavu. Fakticky fungoval jako bezpečný nástroj pro čtení dat ze \textit{Secure Storage}, který prokazoval, zda předchozí operace (např. uložení tokenu po loginu či modifikace preferované brány po setupu) proběhly úspěšně. V pozdějších fázích projektu umožňoval navíc zobrazení aktivní chybové hlášky nebo logu připojení z backendu v reálném čase.

\section{Pokročilé funkce a technické výzvy}
Během integrace stěžejních funkčních celků do architektury se objevilo několik netriviálních technických výzev, které vyžadovaly hlubší analýzu a návrh specializovaných řešení.

\subsection{Správa VPN spojení (OpenVPN)}
Samotné vytvoření šifrovaného tunelu skrze protokol OpenVPN těží z již existující architektury modulu \texttt{ProtoService}, který je nasazen na ostatních platformách. Pro potřeby CLI klienta tak stačilo mapovat požadavky z UDS rozhraní do podoby interních struktur sítě a volat sdílenou aplikační logiku. Modul automaticky vygeneruje odpovídající konfiguraci, sestaví spojení a na pozadí monitoruje stabilitu relace.

\subsection{Persistentní relace a vícero uživatelů}
Jednou z definovaných funkčností byla možnost navázání persistentního (trvalého) spojení, kdy VPN zůstává aktivní i po odhlášení aktuálního systémového uživatele nebo po restartu zařízení.
Tento požadavek s sebou přinesl zásadní problém týkající se izolace více Linuxových uživatelů.

Jelikož je sestavené VPN spojení systémového charakteru (routování přes rozhraní tunelu ovlivňuje globální směrovací tabulky stroje), mohlo by dojít k situaci, kdy jeden uživatel sestaví persistentní spojení a jiný uživatel po přihlášení nemá oprávnění jej modifikovat či ukončit, přestože je jeho síťový provoz rovněž šifrován.
Z toho důvodu bylo nutné obohatit komunikační protokol IPC o explicitní předávání systémového kontextu a identifikaci uživatele, který požadavek vznáší. Pokud je aktivní persistentní tunel spuštěný jiným uživatelem a nejedná se o autorizovanou operaci, je akce zamítnuta na úrovni služby. Administrátor pak musí relaci manuálně ukončit přes příkaz vyžadující superuživatelská práva, např. voláním \texttt{sudo ga-cli disconnect}.

\subsection{Implementace protokolu WireGuard}
Moderní protokol WireGuard představoval v kontextu architektury významnou výzvu, jelikož práce zahrnovala implementaci samotného agenta (sestavování sítě) cíleně pouze pro prostředí Linux s využitím statických konfigurací.
Na rozdíl od uceleného řízení (životní cyklus přes operační systém, dynamické doporučování gateway, automatické znovupřipojení), které přesahuje rámec této bakalářské práce a je paralelně vyvíjeno ostatními členy týmu na různých platformách, se tato konkrétní implementace zaměřuje primárně na fundamentální logiku navázání komunikace přes nativní systémové nástroje (ip link, wg). Tyto operace jsou vloženy do kontextu služeb \texttt{ProtoService} a tvoří základ budoucí robustní síťové logiky.

\subsection{Vylepšení uživatelské přívětivosti (UX)}
S cílem maximalizovat komfort koncových uživatelů při interakci v terminálu byl klient postupně doplňován o podpůrné funkce a takzvané přepínače (flags).
Příkladem je plná podpora pro strojově čitelné formátování výstupů prostřednictvím flagu \texttt{--json}, jenž zajišťuje, že aplikace nebude vracet interaktivní textové uživatelské rozhraní, ale naopak striktně strukturovaný JSON výstup. Tato funkce je naprosto nezbytná pro možnou integraci CLI do DevOps pipelinek a automatizačních skriptů, kde je vyžadováno bezobslužné řízení celého procesu připojení. Další UX optimalizace zahrnovaly vylepšení asynchronního logování stavu a přehlednější zpracování konfiguračních interakcí.
